% (C) 2026 Ralf Stephan, in collaboration with Claude Code.
% Released under CC0 1.0 Universal (public-domain dedication).
% See https://creativecommons.org/publicdomain/zero/1.0/

\documentclass[11pt]{article}
\usepackage[T1]{fontenc}
\usepackage[utf8]{inputenc}
\usepackage{lmodern}
\usepackage[margin=2.5cm]{geometry}
\usepackage{amsmath,amssymb}
\usepackage{braket}
\usepackage[hidelinks]{hyperref}

% Lean names and paths: verbatim, breakable at . _ /
\DeclareUrlCommand\lean{\urlstyle{tt}}
\def\UrlBreaks{\do\.\do\_\do\/}
% Marks an open problem that blocks checking a main claim of the brief
\newcommand{\blocker}[1]{\textsc{[blocks verification: #1]}~}

\title{Open problems:\\ Superpositions correspond to multiple decompositions}
\author{}
\date{}

\begin{document}
\maketitle

\noindent
Open problems raised by the technical brief in this directory, and by its Lean formalization
(see \lean{formalization-log.tex}).
Items marked \textsc{[blocks verification]} stand in the way of checking a main
claim of the brief as the brief states it: a \emph{formalization gap} can be closed
in Lean alone, while a \emph{brief issue} needs the brief to be corrected or made
precise first.

No item below blocks the verification of a main claim of the brief.

\section*{From the brief}
\begin{enumerate}
\item \textbf{No classical superpositions.} The introduction argues that ``mixtures of different
  classical states always yield a different ensemble'', so that classical superpositions are not
  possible. The uniqueness of the decomposition of a classical ensemble (a probability measure)
  into pure (Dirac) states, and the resulting failure of the classical analogue of QSUP-MIX, are
  not formalized. (Brief, Sec.~1.)
\item \textbf{Superpositions of more than two states.} The brief treats superpositions of two
  states. Whether a superposition of $n$ linearly independent states corresponds in the same way
  to an ensemble with a decomposition into those $n$ states and one into $\phi$ and some other
  states is not addressed. (Brief, Sec.~2, QSUP-PUR, QSUP-MIX.)
\end{enumerate}

\section*{From the formalization}
\begin{enumerate}
\item \textbf{Density operators.} QSUP-MIX is stated for convex combinations of pure states
  without the trace; that $\rho$ is a density operator in the sense of QST (positive, trace one,
  and trace class in infinite dimension) is not checked. (Log, ``Differences''.)
\end{enumerate}

\end{document}
